\begin{abstract}
The deployment cost of NMS-free detectors is usually
equated with the latency of the compiled engine. We measure the full
decode-to-detections pipeline of YOLO26n on a Jetson Orin Nano Super:
module input power, clocks, GPU load and official COCO accuracy
on all 5,000 val2017 images. The same FP16 TensorRT engine that takes
4.2\,ms back-to-back takes 13.2\,ms inside a sequential application,
because the stock GPU governor holds the GPU at minimum clock. A
duty-cycle model describes the measurements: engine time scales with
the inverse of the clock (validated by a fixed-frequency sweep to
4\%), interior-clock operating points sit in a 0.53--0.74
GPU-load band while floor-clock and unlocked points fall outside,
and a sequential pipeline released from the top
clock falls to the floor in seconds. Restructuring only the host
pipeline, under default governors and with identical predictions,
gives $2.4\times$ the throughput of the best sequential pipeline at
locked maximum clocks. The CPU governor is coupled to the same
loop: a sleeping CUDA wait lets \texttt{schedutil} lower the host's
clock and costs a fifth of throughput; a per-task utilization
clamp, needing no root privileges, recovers most of it at
lower energy. At camera rates below capacity, idle power sets energy
per frame and the schedule sets latency. Replications on a
Raspberry Pi~5 with Hailo-8 NPU and a Qualcomm QCM6490 (CPU and
Hexagon NPU) show the effect transfers while its mechanism shifts to
the host side. Edge-detection efficiency should be reported for the
whole pipeline under deployment governors.
\end{abstract}
